\subsection{Poisson 公式}

在本节中，我们考察 G 的一个闭子群 H。我们用 \(\beta = dx\) 表示 G 上的 Haar 测度，用 \(\widehat{\beta}\) 表示 \(\widehat{G}\) 上的对偶 Haar 测度。用 \(\alpha\) 表示 H 上的一个 Haar 测度，用 \(\widehat{\alpha}\) 表示对偶群 \(\widehat{H}\) 上的对偶 Haar 测度，我们把后者与 \(\widehat{G}/H^{\perp}\) 等同（II，第

226 页的定理 4）。我们还把 \(\widehat{\mathrm{G}}/\widehat{\mathrm{H}}\) 与 \(\mathrm{H}^{\perp}\) 通过典范投影 \(\mathrm{G}\to\mathrm{G}/\mathrm{H}\) 的对偶映射等同起来（同前）。

我们将用 \(\dot{x}\) 表示元素 \(x\)（属于 G）在 G/H 中的典范像，用 \(\dot{\chi}\) 表示元素 \(\chi\)（属于 \(\widehat{G}\)）在 \(\widehat{G}/H^{\perp}\) 中的典范像。

我们用 \(\gamma\) 表示 G/H 上的 Haar 测度 \(\beta/\alpha\)（INT，VII，§2，第 2 节，定义 1 及第 7 节，命题 10），用 \(\widehat{\gamma}\) 表示 \(H^{\perp}\) 上的对偶 Haar 测度。回顾（INT，VII，§2，第 3 节，命题 5, c)）：测度 \(\gamma\) 由下述性质刻画：对每个 \(f \in \mathcal{L}^{1}(G)\)，H 上的函数 \(y \mapsto f(xy)\) 是 \(\alpha\)-可积的（对 \(\beta\)-几乎所有的 \(x \in G\)）；其积分只依赖于 \(\dot{x}\)，且在 G/H 上由以下方式 \(\gamma\)-几乎处处定义的函数

\[
f ^ {\flat} \colon \dot {x} \mapsto \int_ {\mathrm{H}} f (x h) d \alpha (h)
\]

属于 L \(^{1}\) (G/H, γ)，且满足

\[
\int_ {\mathrm{G} / \mathrm{H}} f ^ {\flat} d \gamma = \int_ {\mathrm{G}} f d \beta .\tag{33}
\]

命题 15.　设 \(f \in \mathrm{L}^1(\mathrm{G})\)，使得在 \(\mathrm{H}^\perp\) 上连续函数 \(\mathcal{F}_{\mathrm{G}}(f)\) 的限制关于 \(\widehat{\gamma}\) 可积。则对几乎所有的 \(x \in \mathrm{G}\)，函数 \(y \mapsto f(xy)\)（它在 \(\mathrm{H}\) 上）是 \(\alpha\)-可积的，且有：

\[
\int_ {\mathrm{H}} f (x y) d \alpha (y) = \int_ {\mathrm{H} ^ {\perp}} \langle \chi , x \rangle \mathcal {F} _ {\mathrm{G}} (f) (\chi) d \widehat {\gamma} (\chi).
\]

根据上述，函数 \(f^{\flat}\)（在 G/H 上几乎处处由

\[
f ^ {\flat} (\dot {x}) = \int_ {\mathrm{H}} f (x y) d \alpha (y)
\]

定义）属于 \(L^{1}(G/H)\)。\(f^{\flat}\) 的 Fourier 变换等同于 \(H^{\perp} = \widehat{G/H}\) 上如下给出的函数：对 \(\chi \in H^{\perp}\)，

\[
\begin{array}{r} \mathcal {F} _ {\mathrm{G/H}} (f ^ {\flat}) (\chi) = \int_ {\mathrm{G/H}} \overline {{\langle \chi , \dot {x} \rangle}} f ^ {\flat} (\dot {x}) d \gamma (\dot {x}) \\ = \int_ {\mathrm{G}} \overline {{\langle \chi , x \rangle}} f (x) d \beta (x) = \mathcal {F} _ {\mathrm{G}} (f) (\chi) \end{array}
\]

这是对可积函数 \(x \mapsto \overline{\langle\chi, x\rangle} f(x)\) 应用公式 (33) 得到的。根据假设，函数 \(\mathcal{F}(f)|\mathrm{H}^{\perp} = \mathcal{F}_{\mathrm{G}/\mathrm{H}}(f^{\flat})\) 属于 \(\mathrm{L}^{1}(\mathrm{H}^{\perp})\)，因而函数 \(f^{\flat}\) 属于空间 B(G/H)。由此（II，第 222 页，定理 3），\(f^{\flat}\) 与 \(\overline{\mathcal{F}}_{\widehat{\mathrm{G / H}}}(\mathcal{F}_{\mathrm{G / H}}(f^{\flat}))\) 几乎处处重合。因此，对几乎每个 \(\dot{x}\in \mathrm{G / H}\)，我们有

\[
f ^ {\flat} (\dot {x}) = \int_ {\mathrm{H} ^ {\perp}} \langle \chi , x \rangle \mathcal {F} _ {\mathrm{G/H}} (f ^ {\flat}) (\chi) d \widehat {\gamma} (\chi) = \int_ {\mathrm{H} ^ {\perp}} \langle \chi , x \rangle \mathcal {F} _ {\mathrm{G}} (f) (\chi) d \widehat {\gamma} (\chi).
\]

这就完成了证明。

推论（Poisson 公式）.　设 \(f \in \mathcal{L}^{1}(G)\)。假设下列条件满足：

\begin{enumerate}
\def\labelenumi{(\roman{enumi})}
\item
  \(\mathcal{F}_{\mathrm{G}}(f)\) 在 \(\mathrm{H}^{\perp}\) 上的限制是可积的；
\item
  对每个 \(x \in G\)，H 上的函数 \(y \mapsto f(xy)\) 是可积的；
\item
  映射 \(x\mapsto \int_{\mathrm{H}}f(xy)d\alpha (y)\) 在 G 上连续。则有
\end{enumerate}

\[
\int_ {\mathrm{H}} f (y) d \alpha (y) = \int_ {\mathrm{H} ^ {\perp}} \mathcal {F} _ {\mathrm{G}} (f) (\chi) d \widehat {\gamma} (\chi).\tag{34}
\]

事实上，沿用上一命题证明中的记号，G/H 上的函数 \(f^{\flat}\) 与 \(\overline{\mathcal{F}}_{\widehat{\mathrm{G/H}}}\left(\mathcal{F}_{\mathrm{G/H}}(f^{\flat})\right)\) 连续且几乎处处相等。因此它们处处相等，特别地在 e 处相等，这就给出公式 (34)。

命题 16.　测度 \(\widehat{\alpha}\)（它在 \(\widehat{\mathrm{H}} = \widehat{\mathrm{G}} / \mathrm{H}^{\perp}\) 上）等于 \(\widehat{\beta} / \widehat{\gamma}\)。

取定一个非零的 \(f\in \mathcal{K}(\mathrm{G})\)。对 \(x\in \mathrm{G}\) 与 \(\chi \in \widehat{\mathrm{G}}\)，令

\[
\varphi (x, \chi) = \int_ {\mathrm{H}} f (x y) \langle \chi , y \rangle d \alpha (y).
\]

函数 \(\varphi\) 在 \(G \times \widehat{G}\) 上连续（INT，IV，§4，第 3 节，定理 2 的推论 1）。固定 \(x\) 时，\(\varphi(x, \chi)\) 只依赖于 \(\chi\) 在 \(\widehat{G}/H^{\perp} = \widehat{H}\) 中的类。固定 \(\chi\) 时，\(\langle \chi, x \rangle \varphi(x, \chi)\) 只依赖于 \(x\) 在 \(G/H\) 中的类，且函数 \(\dot{x} \mapsto \langle \chi, x \rangle \varphi(x, \chi)\)（在 \(G/H\) 上）具有紧支集。

设 \(x \in G\)。函数 \(\dot{\chi} \mapsto \varphi(x, \chi)\)（在 \(\hat{H}\) 上）是 H 上函数 \(y \mapsto f(xy)\) 的 Fourier 余变换。后者是平方可积的，故由 II，第 217 页的 Plancherel 公式 (23)，我们有

\[
\int_ {\widehat {\mathrm{G}} / \mathrm{H} ^ {\perp}} | \varphi (x, \chi) | ^ {2} d \widehat {\alpha} (\dot {\chi}) = \int_ {\mathrm{H}} | f (x y) | ^ {2} d \alpha (y).\tag{35}
\]

 设 \(\chi\in\widehat{\mathrm{G}}\)。函数 \(\dot{x}\mapsto\langle\chi,x\rangle\varphi(x,\chi)\) 属于 \(\mathcal{H}(\mathrm{G}/\mathrm{H})\)，从而属于 \(\mathrm{L}^{1}(\mathrm{G}/\mathrm{H})\)。它的 Fourier 余变换是 \(\mathrm{H}^{\perp}\) 上的函数

它在 \(\xi\in\mathrm{H}^{\perp}\) 处的值为

\[
\begin{array}{r l} & {\int_ {\mathrm{G/H}} \langle \xi , \dot {x} \rangle \langle \chi , x \rangle \varphi (x, \chi) d \gamma (\dot {x}) = \int_ {\mathrm{G/H}} \Bigl (\int_ {\mathrm{H}} \langle \chi \xi , x y \rangle f (x y) d \alpha (y) \Bigr) d \gamma (\dot {x})} \\ & {\qquad = \int_ {\mathrm{G}} \langle \chi \xi , x \rangle f (x) d \beta (x) = \overline {{\mathcal {F}}} _ {\mathrm{G}} (f) (\chi \xi)} \end{array}
\]

（由公式 (33)）。因此

\[
\int_ {\mathrm{G/H}} | \varphi (x, \chi) | ^ {2} d \gamma (\dot {x}) = \int_ {\mathrm{H} ^ {\perp}} | \overline {{\mathcal {F}}} _ {\mathrm{G}} (f) (\chi \xi) | ^ {2} d \widehat {\gamma} (\xi)\tag{36}
\]

这是再次由 Plancherel 公式得到的。

最后我们计算

\[
\begin{array}{l l} \int_ {\widehat {G}} | \overline {{\mathcal {F}}} _ {G} (f) | ^ {2} d \widehat {\beta} = \int_ {G} | f | ^ {2} d \beta & (\text {by (23)}) \\ = \int_ {G / H} d \gamma (\dot {x}) \int_ {H} | f (x y) | ^ {2} d \alpha (y) & (\text {by (33)}) \\ = \int_ {G / H} d \gamma (\dot {x}) \int_ {\widehat {G} / H ^ {\perp}} | \varphi (x, \chi) | ^ {2} d \widehat {\alpha} (\dot {\chi}) & (\text {by (35)}) \\ = \int_ {\widehat {G} / H ^ {\perp}} d \widehat {\alpha} (\dot {\chi}) \int_ {G / H} | \varphi (x, \chi) | ^ {2} d \gamma (\dot {x}) \\ = \int_ {\widehat {G} / H ^ {\perp}} d \widehat {\alpha} (\dot {\chi}) \int_ {H ^ {\perp}} | \overline {{\mathcal {F}}} _ {G} (f) (\chi \xi) | ^ {2} d \widehat {\gamma} (\xi) & (\text {by (36)}) \end{array}
\]

其中把 INT，V，§8，第 3 节，命题 5 应用于连续正函数 \((\dot{x},\dot{\chi})\mapsto |\varphi (x,\chi)|^{2}\)（在 \(\mathrm{G / H}\times \widehat{\mathrm{G}} /\mathrm{H}^{\perp}\) 上）。

把这个等式与群 \(\widehat{G}\) 的积分公式 (33) 相比较，便可断定 Haar 测度 \(\widehat{\alpha}\) 与 \(\widehat{\beta}/\widehat{\gamma}\) 重合。

